\chapter{Introduction to Numerical Methods for Optimization}

\section{Motivations and introductory example}

\emph{Optimization} consists of choosing, among a set of possible decisions,
the one that minimizes (or maximizes) a quantity called the \emph{objective function}.
In this course, we mainly focus on \emph{minimization} problems, although
maximization problems can be treated equivalently by considering the opposite of the function.

\medskip

Optimization problems encountered in practice can be:

\begin{itemize}
  \item \textbf{continuous}, when variables take values in $\R^n$;
  \item \textbf{discrete}, when some variables must be integers
        (integer programming);
  \item \textbf{mixed} (MIP: Mixed Integer Programming), when some variables
        are continuous and others are discrete.
\end{itemize}

Discrete or mixed problems are generally harder to solve than continuous
problems and will not be studied in detail in this course.

\medskip

We also encounter problems where parameters are random. This leads to:

\begin{itemize}
  \item \textbf{stochastic optimization}, where we minimize an expectation;
  \item \textbf{chance constraints} (chance-constrained), where certain constraints must
        be satisfied with high probability;
  \item \textbf{robust optimization}, where constraints must be satisfied for all
        admissible realizations of an uncertainty set.
\end{itemize}

These frameworks are essential in fields such as finance, energy management
or systems engineering, but they fall outside the immediate scope of this chapter.

\begin{exemple}[Quadratic problem with linear constraints]
We consider the problem
\begin{equation}
  \min_{x \in \R^2} \; f(x)
  := (x_1 - 1)^2 + 2(x_2 + 3)^2
  \label{eq:intro-exemple-obj}
\end{equation}
under the constraints
\begin{equation}
  x_1 \geq 2, 
  \qquad
  x_1 - x_2 = 4.
  \label{eq:intro-exemple-constraints}
\end{equation}
\end{exemple}

The objective function $f$ measures the ``badness'' of a vector $x=(x_1,x_2)$.
The smaller $f(x)$ is, the more satisfactory the solution is. The constraints impose that the
admissible solutions live on a line and in a half-plane. The level curves
of $f$ are ellipses centered at $(1,-3)$.

\begin{figure}[H]
  \centering
  \begin{tikzpicture}[scale=0.9]
    % Hatched region for x >= 2 (feasible zone)
    \fill[pattern=north east lines, pattern color=gray!30] (2,-6) rectangle (5,1);
    
    % Axes - thicker for clarity
    \draw[->, very thick] (-0.5,0) -- (5,0) node[right] {$x$};
    \draw[->, very thick] (0,-6) -- (0,1) node[above] {$y$};

    % Center
    \fill (1,-3) circle (1.5pt) node[below right] {$(1,-3)$};

    % Levels (approximate ellipses)
    \draw[blue, thick] (1,-3) ellipse (1.0 and 0.7);
    \draw[blue]        (1,-3) ellipse (2.0 and 1.4);
    \draw[blue]        (1,-3) ellipse (3.0 and 2.1);

    % Half-plane x >= 2
    \draw[gray, thick] (2,-6) -- (2,1) node[above] {$x = 2$};
    \node[gray] at (4.2,-5.5) {$x \ge 2$};

    % Line x - y = 4  => y = x - 4
    \draw[red, thick, domain=-0.5:5] plot (\x, {\x - 4}) node[above right] {$x - y = 4$};

    % Feasible solution (intersection constraint + low level)
    \coordinate (xmin) at (2,-2);
    \fill[red] (xmin) circle (2pt) node[above left] {$x^\star$};
  \end{tikzpicture}
  \caption{Illustration of the introductory quadratic problem.}
  \label{fig:intro-exemple}
\end{figure}

\begin{figureexplanation}
In Figure~\ref{fig:intro-exemple}, the \textbf{blue ellipses} represent the level curves of the objective function~$f$: the smaller the ellipse, the lower the value of~$f$. The \textbf{center} of the ellipses, located at $(1,-3)$, corresponds to the unconstrained minimum of~$f$. The \textbf{hatched area} represents the feasible region $x \ge 2$. The \textbf{vertical gray line} $x = 2$ materializes the boundary of this constraint, and the \textbf{red line} represents the equality constraint $x - y = 4$. The point $x^\star$, intersection of the red line with the boundary of the feasible region, is the \textbf{optimal solution} of the constrained problem.
\end{figureexplanation}

This problem illustrates three essential components:
\begin{itemize}
  \item the \textbf{objective function} $f : \R^n \to \R$;
  \item the \textbf{set of constraints}, defining the feasible points;
  \item the \textbf{optimal solution} $x^\star$, minimizing $f$ among the feasible points.
\end{itemize}

\section{General formulation of an optimization problem}

\begin{definition}[Optimization problem]
An \emph{optimization problem} is the search for a vector
$x \in \R^n$ minimizing a function $f : \R^n \to \R$ under equality
and/or inequality constraints, that is:
\begin{equation}
  \begin{aligned}
    \min_{x \in \R^n} \quad & f(x) \\
    \text{under the constraints} \quad &
    c_i(x) \ge 0,\quad i \in I, \\
    & c_j(x) = 0,\quad j \in E.
  \end{aligned}
  \label{eq:general-problem}
\end{equation}
\end{definition}

\begin{definition}[Feasible set]
The feasible set is
\[
\mathcal{X} = \bigl\{x \in \R^n : c_i(x)\ge 0\ \forall i\in I,\; c_j(x)=0\ \forall j\in E\bigr\}.
\]
Any point in $\mathcal{X}$ is called \emph{feasible} (or admissible).
\end{definition}

\begin{definition}[Optimal solution]
A point $x^\star \in \mathcal{X}$ is a \emph{global optimal solution} if
\[
f(x^\star) \le f(x), \qquad \forall x \in \mathcal{X}.
\]
We speak of a \emph{local minimum} if this property only holds in a neighborhood of $x^\star$.
\end{definition}

\begin{remarque}
In practice, it is rare to determine $x^\star$ exactly.
Numerical methods produce a sequence $(x_k)$ whose gradient values
and/or function values decrease towards zero.
\end{remarque}

\section{Classification of optimization problems}

\subsection{Unconstrained or constrained}

\begin{itemize}
  \item \textbf{Unconstrained optimization}:  
  \[
  \min_{x\in\R^n} f(x).
  \]

  \item \textbf{Constrained optimization}:  
  \[
  x \in \mathcal{X} \subsetneq \R^n.
  \]
\end{itemize}

A common special case is that of box constraints:
\[
\ell_i \le x_i \le u_i, \qquad i=1,\dots,n.
\]

\subsection{Continuous, discrete and stochastic optimization}

\paragraph{Integer linear programming.}
When certain variables $x_i$ must be integers, the problem becomes
\emph{discrete} (NP-hard in general).

\paragraph{Mixed problems (MIP).}
Combination of continuous and discrete variables.

\paragraph{Stochastic optimization.}
Presence of random parameters:
\begin{itemize}
  \item minimization of an expectation;
  \item constraints satisfied with probability $\ge p$ (chance-constrained);
  \item satisfaction for all possible scenarios (robust optimization).
\end{itemize}

\subsection{Linear, quadratic, nonlinear problems}

\begin{itemize}
  \item \textbf{Linear Programming (LP)}:
  \[
  f(x)=c^\top x,\qquad c_i(x)=a_i^\top x - b_i.
  \]

  \item \textbf{Quadratic Programming (QP)}:
  \[
  f(x)=\tfrac12 x^\top Q x + c^\top x,
  \]
  where $Q$ is symmetric.

  \item \textbf{Nonlinear Programming (NLP)}:
  $f$ and/or the constraints are nonlinear.
\end{itemize}

\subsection{Convexity, concavity and convex programs}

\begin{definition}[Convex set]
A set $\mathcal{C}$ is convex if  
\[
\theta x + (1-\theta)y \in \mathcal{C},\qquad \forall x,y\in\mathcal{C},\;\theta\in[0,1].
\]
\end{definition}

\begin{definition}[Convex function]
A function $f$ is convex if  
\[
f(\theta x + (1-\theta)y)\le \theta f(x)+(1-\theta)f(y).
\]
It is \emph{strictly convex} if the inequality is strict for $x\neq y$.
It is \emph{concave} if $-f$ is convex.
\end{definition}

\begin{definition}[Convex program]
An optimization problem is called \emph{convex} when:
\begin{itemize}
  \item the objective function $f$ is convex;
  \item the equality constraints are linear;
  \item the inequality constraints $c_i(x)\ge 0$ are concave.
\end{itemize}
\end{definition}

In this framework, any local minimum is automatically a global minimum.

\subsection{Penalty methods}

Certain constraints can be integrated into the objective function.  
For example:
\[
\min_{x\in\R^2} (x_1-1)^2+(x_2+3)^2\quad\text{subject to}\quad x_1+x_2=5
\]
can be reformulated as an unconstrained problem:
\[
\min_{x\in\R^2} (x_1-1)^2+(x_2+3)^2
+ \frac{\eta}{2}\,(x_1+x_2-5)^2, \qquad \eta>0.
\]

We will discuss this again in the last chapter on penalization, but this second problem is only equivalent to the first for $\eta\to\infty$. Precisely, its solution is 
\[
x_1 = \frac{9\eta +2}{2\eta +2}, \quad x_2 = \frac{\eta - 6}{2\eta + 2}.
\]

Alternatively, we can consider the problem 
\[
\min (x_1-1)^2 + (x_2+3)^2 +\eta |x_1+x_2 -5|,
\] 
for which we will indicate in the last chapter that the solution coincides with that of the initial constrained problem as soon as $\eta > 7$.

\section{General principle of optimization algorithms}

Numerical optimization methods are \emph{iterative}: they build a sequence
$(x_k)$ converging towards an optimal (or stationary) solution. The general scheme is:
\[
x_{k+1}=x_k+\alpha_k p_k,
\]
where $p_k$ is a search direction and $\alpha_k>0$ is a step length.

\begin{center}
\fbox{
\begin{minipage}{0.9\textwidth}
\textbf{Input:} an initial point $x_0$.\\[0.2em]
\textbf{For} $k=0,1,2,\dots$ \textbf{while} the stopping criterion is not satisfied:
\begin{enumerate}
  \item Choose a direction $p_k$ (for example $p_k=-\nabla f(x_k)$).
  \item Choose a step length $\alpha_k>0$ (line search, trust region, etc.).
  \item Update $x_{k+1}=x_k+\alpha_k p_k$.
\end{enumerate}
\textbf{End For}
\end{minipage}}
\end{center}

For unconstrained methods, the algorithm relies on the following information:

\begin{itemize}
  \item only $f(x_k)$ (``black box'' methods);
  \item the gradient $\nabla f(x_k)$ (gradient methods);
  \item the Hessian $\nabla^2 f(x_k)$ (Newton-type methods);
  \item approximations of the Hessian (quasi-Newton methods).
\end{itemize}

\section{Desirable properties of numerical methods}

An efficient optimization method must satisfy several criteria:

\begin{itemize}
  \item \textbf{Robustness}: good behavior for different problems and starting points;
  \item \textbf{Efficiency}: reasonable cost in time and memory;
  \item \textbf{Accuracy}: low sensitivity to rounding errors and calculation inaccuracies;
  \item \textbf{Numerical stability}: avoiding the amplification of errors;
  \item \textbf{Reliability}: avoiding stagnations or chaotic behaviors.
\end{itemize}

\section{Stopping criteria}

An optimization algorithm generally stops when one of the following criteria is reached:

\begin{itemize}
  \item \textbf{small gradient norm}:
        \[
        \|\nabla f(x_k)\| \le \varepsilon_{\mathrm{grad}};
        \]
  \item \textbf{small variation in the solution}:
        \[
        \|x_{k+1}-x_k\|\le \varepsilon_{\mathrm{step}};
        \]
  \item \textbf{small variation in the function}:
        \[
        |f(x_{k+1})-f(x_k)|\le \varepsilon_{\mathrm{f}};
        \]
  \item \textbf{maximum number of iterations reached}.
\end{itemize}

\medskip

This first chapter establishes the general framework of numerical optimization: formulation, classification,
convexity, broad typology of problems, algorithm structure, and essential properties.
The following chapters will be devoted to the detailed study of solving methods,
\bigskip
\begin{checkpoint}
\begin{itemize}
    \item \textbf{Vocabulary:} Can you define: objective function, feasible set, local vs global solution?
    \item \textbf{Structure:} What are the ingredients of a descent algorithm (direction + step)?
    \item \textbf{Convexity:} Why is it important? (A local min is global).
    \item \textbf{Stopping criteria:} Cite two conditions to stop an algorithm.
    \item \textbf{Positive definite Hessian:} $\nabla^2 f(x^\star)$ is positive definite if $\eta^\top \nabla^2 f(x^\star) \eta > 0$ for all vectors $\eta \neq 0$.
\end{itemize}
\end{checkpoint}
beginning with unconstrained optimization.
